\chapter{Celková koncepce řešení}

\section{Game Manager}
Videohry se obvykle skládají z různých scén, které mohou představovat jednotlivé levely, tutorial část, hlavní menu nebo závěrečné titulky. Kdyby videohry musely načítat všechny tyto scény najednou, zabralo by to zbytečně moc času a výpočetní síly i přesto, že se hráč nachází jen na jednom z těchto míst. To je jeden z příkladů optimalizace a logiky, které jsou klíčové pro plynulý chod videohry. Pro tyto účely je jsou používány objekty typu \emph{Game Manager} (dále také „herní manažer“).

Herních manažerů se ve videoherních projektech vyskytuje hned několik a každý má na starosti odlišnou funkční část videohry: Game Manager (základ -- stará se o přepínání mezi scénami, správné zmrazení času, jednoduché uživatelské rozhraní…), Audio Manager (zajišťuje správný chod veškeré logiky zvuku   v prostředí hry nebo v UI), Cursor Manager (upravuje vlastnosti kurzoru počítače pro účely videohry) a další manažery, které nejsou zas tak časté, ale mohou se hodit pro konkrétní případy (Tutorial Manager, Collectible Manager, DontDestroyOnLoad, Death Manager).

Herní manažer je (zejména v Unity) tedy „empty GameObject“ ve scéně, který má na sobě připojený stejnojmenný skript jako komponent, do kterého můžeme přidávat reference k dalším objektům ve scéně nebo v projektu (tělo hráče, další scéna, textové pole) a upravovat jakékoli parametry tohoto skriptu, které zpřístupníme Unity Editoru.

\section{Umírání}
Videohra je navrnuta tak, že každý level má dvě části: \emph{The Chase} a \emph{The Exploration}. V \emph{The Chase} je hráč naháněn masivní bowlingovou koulí, která, pokud hráče dožene a přejede ho, tím ho i zabije a hráč bude muset začít znovu. Toto je způsob smrti \emph{hard death}, kdy se po zabití hráče znovu načte celý level, všechny parametry se restartují a hráč si tak musí \emph{The Chase} zapamatovat celou nazpaměť včetně všech překážek.

\emph{The Exploration} využívá více milosrdnou \emph{soft death}, kdy, pokud hráč spadne nebo se mu nějak jinak povede zemřít, místo toho, aby se scéna načetla od začátku, tak je hráč jen přemístěn k poslednímu aktivovanému checkpointu. Lze říct, že tato videohra tak používá hybridní systém smrti.

\subsection{Checkpoint System}
Do této videohry je dobrý nápad zakomponovat \emph{Checkpoint System}, který zajistí, že pokud hráč bude umírat v pozdějších fázích daného levelu, nemusí celým levelem procházet znova, ale hra ho transportuje k nejbližšímu checkpointu, který si uložil.

Některé hry používají i auto-save systém, kdy hra ukládá progres skoro každým momentem, tudíž kdykoli videohru můžete vypnout, zapnout a objevit se na tom samém místě, kde jste hru opustili, avšak to je ideálnější spíš pro videohry žánru open-world a v této videohře bude stačit, pokud hráč bude docházet k checkpointům a aktivovat je.

\section{Emotion System}
Pro zpestření herního zážitku této videohry byl vymyšlen koncept \emph{Emotion System}. Hráč bude muset v části \emph{The Exploration} sbírat \emph{Conversation~Fragments} (úlomky konverzace), které obsahují atributy emocí, a na základě kontextu celé konverzace bude hráč muset správně určit, jestli má daný \emph{Conversation~Fragment} pozitivní, negativní, nebo neutrální emoci. Podle toho pak bude potřeba sesbírané \emph{Conversation~Fragments} správně umístit na \emph{The Altar} (oltář), který je schopen aktivovat portál do dalšího levelu v případě správné kombinace.

Díky tomuto konceptu videohra bude obsahovat i mírný puzzle-solving a integraci příběhu, které posílí atmosféru surrealismu, kterého videohra má docílit. Je totiž důležité, aby videohra neobsahovala pouze akci, ale i části, kde si hráč může oddechnout a připravit se na další akci. Stejně jako v hudbě, kde je komplikované mít majestátní, hlasité momenty, pokud není správně zacházeno s tichem. Vše je o kontrastu.

\section{UI}
Velkou kapitolou, která musí být v nějakém měřítku v každé videohře, je uživatelské rozhraní (dále také anglická zkratka UI). Jakákoli tlačítka, nastavení herního zážitku nebo interakce se světem videohry musí projít přes skripty, které řeší uživatelské rozhraní.

Jedna z nejzákladnějších částí UI je Input System, který řeší, jaké hardware vstupy způsobují jakou akci ve videohře. Dvě základní vstupní zařízení, kterými dnešní počítače disponují, jsou klávesnice a myš. Některé videohry používají pouze klávesnici (Hollow Knight od studia Team Cherry \cite{hollow-knight}) a některé zas pouze myš (Happy Game od studia Amanita Design \cite{happy-game}). Častěji se však můžete setkat s kombinací obou. Dříve byla ve vývoji této videohry používána kombinace obou -- levá ruka obsluhovala veškerý pohyb včetně skákání a dashe (rychlý posun kupředu), zatímco pravá ruka zaujímala pozici na myši, která byla využívána pro míření, střílení a „camera peeking". Postupem času se v projektu však přešlo výhradně na klávesnici, jelikož tato varianta lépe rozloží aktivitu mezi obě ruce a eliminuje potřebu přesunu ruky mezi zařízeními. V budoucnu je zvažována implementace více možností ovládání, aby si každý hráč mohl vybrat variantu, která mu nejvíce vyhovuje.

Poté, co videohra ví, jaké vstupy mají jaký učel, je třeba s těmito signály pracovat a používat je pro kontroly, mechaniky a orientaci v programu videohry. Nejčastější případ UI jsou panely, které mají různé účely. Například Pause Panel je lišta, která se objeví při pozastavení hry (nejčastěji stisknutím klávesy Escape), kde se mohou objevovat informace o aktuálním stavu hráče, nastavení hlasitosti zvuku nebo možnosti opuštění levelu a ukončení hry.

Dokonce hned první věc, kterou hráč po zapnutí videohry vidí, je nejčastěji obrazovka hlavního menu, která se ve většině případech skládá pouze z UI komponentů. Častá jsou tlačítka „Hrát“, „Nastavení“, „Ukončit“ a „Titulky“ s pozadím, názvem herního titulu a Main Menu hudbou hrající na pozadí.

\section{Vizuální efekty}
Vizuální efekty by, jako každá část, měly podpořit uvěřitelnost digitálního světa a angažovanost do něj. Neměly by na sebe lákat až moc pozornosti, pokud tak vyloženě nechceme. Příklady klasických vizuálních efektů ve hrách jsou:

\begin{itemize}
    \item efekt smrti (Death Screen) 
    \item impact effect (například dopad na zem, dash)
    \item glow effect (na hráči, aby viděl kolem sebe, nebo na jiných předmětech, aby byly rychle zaregistrovatelné)
\end{itemize}

Dosáhnout těchto efektů je možno různými způsoby: přidání prvků světla na určitý GameObject \cite{unity-point-light}, experimentování s Unity Particle System \cite{unity-particle-system} nebo manipulování určitývh prvků Post-processing.

\newpage

\section{Audio}

\subsection{OST}

Skládat hudbu pro hru je samozřejmě více kreativní postup, kdy její tvůrce může strávit hodiny sezením u hudebního nástroje nebo vybíráním či tvorbou ideálních zvuků, které má v hlavě. Hudba většinou podporuje již hotovou vizuální a příběhovou stránku videoher (nebo i filmů) -- málokdy se stává, že se svět videohry staví kolem hudby. Tato videohra má surreální, snovou a tajemnou atmosféru, jelikož sleduje příběh postavy, která hledá sebe sama, má neustálý strach, neví, co může čekat od nastávající chvíle, a chce se dostat do bezpečí. Během této cesty se potkává s nadpřirozenými jevy, které nedávají smysl a je v neustálém stavu zmatení.

Pro tato témata byly vybrány následující sonické elementy:

\begin{itemize}
    \item vzdušné, retro tóny inspirované interpretem Joe Hisaishim (populární japonský skladatel pro legendární filmy od Studio Ghibli) \cite{ghibli}, zejména skladbou „The Path of the Wind – Instrumental“\cite{hisaishi-path-wind}
    \item podivné reverb \& delay efekty
    \item mixolydický modus, dórský modus, ukrajinská dórská / rumunská mollová stupnice
    \item mikrotonální hudba, xenharmonie
    \item vinylový efekt
    \item Další inspirace byla převzata například ze skladeb „It’s Just a Burning Memory“ (The Caretaker) \cite{caretaker-burning-memory}, „Moog City 2“ (C418) \cite{c418-moog-city-2} nebo „Resonance“ (Home) \cite{home-resonance}.
\end{itemize}

\subsection{SFX}
SFX (sound effects), neboli zvukové efekty, jsou dnes brány už jako samozřejmost a lidé si tolik neuvědomují, kolik práce na pozadí odvádí. Tvorba zvukových efektů do videohry je tedy další velká část kreativní činnosti, kde je důležité představit si, jak celý svět zní, a vytvořit příslušný zvukový efekt pro každou věc, která by v reálném světě vydávala nějaký zvuk. Sem spadají ambience na pozadí (padání kapek v jeskyních, šumění listů ve vánku, šeptání lidí), zvuky konkrétních objektů / akcí (kroky, útok, aktivace portálu, sbírání předmětů, dopad na zem, simulace komunikace s ostatními postavami) a zvuky uživatelského rozhraní (zmáčknutí tlačítka, pozastavení hry, hover akce).

Zvukové efekty je tak možno rozdělit do dvou částí: diegetické a nediegetické. Diegetické jsou součástí prostředí hry, mohou je slyšet postavy ve hře, zatímco nediegetické jsou určené pro hráče, nespadají do prostředí hry a mají čistě funkční účely (zpětná vazba, indikátory…). V některých případech se v médiích používají i přechody mezi nediegetickými a diegetickými zvuky, které také mohou mít svůj význam. Můj oblíbený případ tohoto efektu je z televizního seriálu Arcane, kde je v první půlminutě této video ukázky slyšet, jak se hudba vycházející z balónu (prostředí světa) postupně přemění na podkladovou hudbu pro diváky.\footnote{Arcane: Season 2 | Ekko and Jinx Team Up (Nerd Clips HD). \url{https://youtu.be/B_KSH6w-gps}}

\section{Post-processing}
Post-processing dodává videohře atmosféru, šmrnc a oživuje její již zhotovenou grafickou stránku. Pracuje se světlem, stínováním, barvami, zrcadlením, šumem a dalšími efekty, které najdete ve většině grafických nástrojů. Tyto efekty jsou pak aplikovány jako poslední grafická vrstva, než obraz dojde na obrazovku, a mohou být zpřístupněny přes skripty a ovládány tak, aby reagovaly na události ve videohře. Mezi časté použití post-processingu patří:

\begin{itemize}
    \item práce s jasem (iluze fáze noci, realistické, intenzivní zesvětlení prostředí simulující explozi nukleární bomby…)
    \item saturace / desaturace obrazu pro vyzdvihnutí emocí
    \item VHS efekt pro videohry žánru „analogový horor“
\end{itemize}